\subsection{19.93: unbounded lower p-central width without a wreath quotient}
\label{cand:19.93}
\problemstatement{19.93}
Write \(P_1(Q)=Q\), \(P_{i+1}(Q)=P_i(Q)^p[P_i(Q),Q]\).
For every fixed prime \(p\ge7\) and integer \(K\), we construct a
finite two-generator \(p\)-group with no quotient \(C_p\wr C_p\)
and some factor \(P_i/P_{i+1}\) of order greater than \(p^K\).

First write \(C_p\wr C_p=\F_p^p\rtimes\langle t\rangle\) and let
\(s(v)\) be the sum of the coordinates of \(v\).
The map \((v,t^a)\mapsto(s(v),a)\) is onto \(\F_p^2\).
The element \((v,t^a)\) has \(p\)-th power one exactly when
\(s(v)a=0\): for \(a\ne0\) its \(p\)-th power has constant vector
\(s(v)\), and for \(a=0\) it is always one.
If \(u,v\) generate the wreath product and \(u^p=v^p=1\), their
images must lie on different coordinate axes, so \((uv)^p\ne1\).
Consequently the three relations
\[
 x^p=y^p=(xy)^p=1                                  \tag{\(\dagger\)}
\]
exclude this quotient for a group generated by \(x,y\).

Consider the complete filtered algebra
\[
 A=\F_p\langle\!\langle X,Y\rangle\!\rangle/
 \overline{(X^p,Y^p,(X+Y+XY)^p)}.
\]
Let \(A_n\) be its degree filtration, \(a_n=\dim A_n/A_{n+1}\),
and \(H(t)=\sum a_nt^n\).
All three relations have initial degree \(p\); in the last relation
the degree-\(p\) part \((X+Y)^p\) is the sum of all distinct
length-\(p\) words and is nonzero.
The complete filtered Golod--Shafarevich inequality
\cite[Theorem~2.6]{ershov-gs} states coefficientwise that
\[
 (1-2t+3t^p)\frac{H(t)}{1-t}\ge\frac1{1-t}.
\]
At \(t=2/3\) the polynomial factor is at most \(-115/729<0\),
so \(H(2/3)\) diverges.
The denominator \(1-t\) is retained: these relations are not homogeneous.

For \(N\ge2\), let \(R_N=A/A_N\) and \(M_N=A_1/A_N\).
The finite group \(1+M_N\) is a \(p\)-group, because \(M_N\) is
nilpotent and \(|1+M_N|=|M_N|\). Set
\[
 Q_N=\langle1+X,1+Y\rangle\le1+M_N.
\]
The defining relations give (\(\dagger\)), excluding the wreath
quotient. Reduction to \(A_1/A_2\cong\F_p^2\) shows that \(Q_N\)
requires exactly two generators.
The natural map \(\F_p[Q_N]\to R_N\) is onto: its image contains
\(X,Y\), and every truncated series is a polynomial.
For its augmentation ideal \(I_N\), the map kills \(I_N^N\);
hence
\[
 \dim R_N\le\dim\F_p[Q_N]/I_N^N.                    \tag{\(\ddagger\)}
\]

Here is an elementary bound on the right side if all lower
\(p\)-central widths were bounded.
For any finite \(p\)-group \(Q\), put
\(b_i=\dim_{\F_p}P_i/P_{i+1}\), and choose basis lifts \(g_{ij}\),
ordered by increasing \(i\). The series terminates, and successive
quotient normal forms give unique products
\(\prod g_{ij}^{e_{ij}}\), \(0\le e_{ij}<p\).
Expanding in this fixed order shows that
\(\prod(g_{ij}-1)^{e_{ij}}\) span \(\F_p[Q]\).
For the augmentation ideal \(I\),
\[
 g-1\in I^i\quad(g\in P_i).
\]
This follows inductively from \(h^p-1=(h-1)^p\) and
\[
 [h,g]-1=h^{-1}g^{-1}
 \bigl((h-1)(g-1)-(g-1)(h-1)\bigr).
\]
Thus only products with \(\sum i e_{ij}<N\) are needed modulo \(I^N\).
If every \(b_i\le K\), and \(c_n\) are the coefficients of
\[
 P_K(t)=\prod_{i\ge1}(1+t^i+\cdots+t^{(p-1)i})^K,
\]
then \(\dim\F_p[Q]/I^N\le\sum_{n<N}c_n\).
This product converges for \(0<t<1\), being bounded by
\(\prod_{i\ge1}(1-t^i)^{-K}\); the logarithms are bounded by
\(K\sum_i t^i/(1-t)\).

If every \(Q_N\) had every \(b_i\le K\), (\(\ddagger\)) would give
\(\sum_{n<N}a_n\le\sum_{n<N}c_n\) for all \(N\).
In particular \(a_n\le\sum_{j\le n}c_j\), and therefore
\[
 H(t)\le\frac{P_K(t)}{1-t}<\infty\qquad(0<t<1),
\]
contradicting divergence at \(2/3\).
This proves the claimed unboundedness, already for \((p,d)=(7,2)\).
The finite groups are obtainable from algebra truncations, but the
proof gives no practical size bound. It excludes wreath quotients
of \(Q_N\) itself, exactly as asked, not quotients of every subgroup.
See \artifact{research/19.93-proof.md}.
